Al utilizar protocolos basados en teoría de grupos, como Diffie-Hellman, ElGamal o firmas de Schnorr, nos basamos en un grupo $G$, un elemento $g$ del grupo y un número natural $q$ que son definidos de antemano por alguna entidad externa. Explique qué son estos elementos $g$ y $q$, cuál es la relación entre ellos. Explique también qué debemos verificar para asegurar que la entidad externa no nos está \emph{engañando} con respecto a esa relación. Justifique por qué dicha verificación asegura que la relación entre $g$ y $q$ es la esperada.
